% ============================================================
% 1.3. Расчётная схема
% ============================================================
\subsection{Расчётная схема}

\subsubsection{Идеализация конструкций}

Переход от реальной конструкции к расчётной схеме является первым и принципиальным шагом инженерного анализа~\cite{feodosev, timoshenko}. В рамках платформы SopromatLab реальные элементы конструкций заменяются идеализированными расчётными схемами на основе следующих допущений~\cite{aleksandrov, belyaev}:

\begin{enumerate}
  \item Конструктивный элемент моделируется как прямолинейный стержень (балка) с осью, совпадающей с осью $x$.
  \item Поперечное сечение балки остаётся постоянным по длине (призматический стержень).
  \item Нагрузки действуют в плоскости $xOy$, проходящей через ось балки (плоская задача изгиба).
  \item Собственный вес балки не учитывается (может быть задан как распределённая нагрузка при необходимости).
\end{enumerate}

\subsubsection{Схематизация нагрузок}

Внешние воздействия на конструкцию схематизируются в следующие идеализированные типы~\cite{darkov, stepin}:

\begin{itemize}
  \item \textbf{Сосредоточенная сила $F$} --- сила, приложенная к бесконечно малому участку оси балки. На расчётной схеме изображается стрелкой с указанием величины и точки приложения.
  \item \textbf{Распределённая нагрузка $q(x)$} --- интенсивность силового воздействия, распределённого по длине. Изображается прямоугольником (для равномерной) или треугольником (для линейно-меняющейся нагрузки).
  \item \textbf{Сосредоточенный момент $M_0$} --- пара сил, приложенная в точке. Изображается дуговой стрелкой.
\end{itemize}

\subsubsection{Граничные условия}

Платформа реализует два типа граничных условий, соответствующих основным видам опор~\cite{feodosev, gere2018}:

\textbf{1. Шарнирная (неподвижная) опора} запрещает вертикальное перемещение в точке закрепления, но допускает поворот:
\begin{equation}
  v(x_{\text{оп}}) = 0, \qquad \theta(x_{\text{оп}}) \neq 0.
  \label{eq:hinge_bc}
\end{equation}

На расчётной схеме изображается треугольником. Реакция --- вертикальная сила $R$.

\textbf{2. Жёсткая заделка (защемление)} запрещает как вертикальное перемещение, так и поворот:
\begin{equation}
  v(x_{\text{зад}}) = 0, \qquad \theta(x_{\text{зад}}) = 0.
  \label{eq:clamp_bc}
\end{equation}

На расчётной схеме изображается штриховкой. Реакции --- вертикальная сила $R$ и реактивный момент $M_{\text{зад}}$.

\subsubsection{Принятые допущения}

Математическая модель базируется на классических гипотезах теории упругости~\cite{timoshenko, timoshenko1983}:

\textbf{1. Гипотеза плоских сечений (гипотеза Бернулли).} Поперечные сечения балки, плоские и перпендикулярные оси до деформации, остаются плоскими и перпендикулярными деформированной оси после деформации:
\begin{equation}
  \varepsilon(x, y) = -y \cdot \kappa(x) = -y \cdot \frac{d^2 v}{dx^2},
  \label{eq:bernoulli}
\end{equation}
\begin{conditions}
  \varepsilon & относительная продольная деформация; \\
  y & расстояние от нейтральной оси; \\
  \kappa(x) & кривизна деформированной оси.
\end{conditions}

\textbf{2. Закон Гука (линейная упругость).} Напряжения пропорциональны деформациям:
\begin{equation}
  \sigma(x, y) = E \cdot \varepsilon(x, y) = -E \cdot y \cdot \frac{d^2 v}{dx^2}.
  \label{eq:hooke}
\end{equation}

\textbf{3. Гипотеза малых перемещений.} Прогибы $v(x)$ малы по сравнению с длиной балки ($v \ll L$), что позволяет заменить кривизну второй производной:
\begin{equation}
  \kappa = \frac{d^2 v / dx^2}{\left[1 + (dv/dx)^2\right]^{3/2}} \approx \frac{d^2 v}{dx^2}.
  \label{eq:small_disp}
\end{equation}

\textbf{4. Отсутствие начальных напряжений.} В недеформированном состоянии напряжения в балке равны нулю.

\textbf{5. Однородность и изотропность материала.} Механические свойства одинаковы во всех точках и по всем направлениям.

Данные допущения обеспечивают линейность задачи и позволяют применять принцип суперпозиции --- складывать результаты от отдельных нагрузок~\cite{aleksandrov, beer2020}.

\subsubsection{Координатные оси и правила знаков}

В платформе SopromatLab приняты следующие правила знаков~\cite{feodosev, belyaev}, реализованные в расчётном ядре:

\begin{itemize}
  \item Ось $x$ направлена вдоль оси балки слева направо (от опоры $A$ к опоре $B$).
  \item Ось $y$ направлена вертикально вверх.
  \item \textbf{Поперечная сила $Q > 0$}, если сумма проекций внешних сил на ось $y$ слева от сечения направлена вверх.
  \item \textbf{Изгибающий момент $M > 0$}, если нижнее волокно балки растянуто (<<балка улыбается>>).
  \item \textbf{Продольная сила $N > 0$} при растяжении стержня.
\end{itemize}

Для визуальной ясности на эпюрах реализована цветовая дифференциация: красный ($\text{\#ef4444}$) --- зоны растяжения ($M > 0$), синий ($\text{\#3b82f6}$) --- зоны сжатия ($M < 0$)~\cite{bostock2011}.

\subsubsection{Вывод формулы нормальных напряжений при изгибе}

Ключевым результатом теории плоского изгиба балок является формула нормальных напряжений, связывающая изгибающий момент в сечении с распределением напряжений по высоте сечения. Приведём её строгий вывод из общих теоретических предпосылок~\cite{feodosev, timoshenko, gere2018}.

\textbf{Шаг~1. Геометрическое соотношение (гипотеза Бернулли).} Рассмотрим элемент балки длиной $dx$. Согласно гипотезе плоских сечений (\ref{eq:bernoulli}), два соседних сечения поворачиваются друг относительно друга на малый угол $d\varphi$. Волокно, расположенное на расстоянии $y$ от нейтральной оси (оси, проходящей через центр тяжести сечения), изменяет свою длину на величину $\Delta(dx) = -y \cdot d\varphi$. Относительная продольная деформация этого волокна:
\begin{equation}
  \varepsilon(y) = \frac{\Delta(dx)}{dx} = -y \cdot \frac{d\varphi}{dx} = -y \cdot \kappa,
  \label{eq:strain_bending}
\end{equation}
где $\kappa = d\varphi / dx$ --- кривизна деформированной оси балки. Знак <<минус>> указывает, что при $\kappa > 0$ (балка изогнута вогнутостью вверх) верхние волокна ($y > 0$) сжимаются, нижние --- растягиваются.

\textbf{Шаг~2. Физическое соотношение (закон Гука).} По закону Гука (\ref{eq:hooke}), нормальные напряжения пропорциональны деформациям:
\begin{equation}
  \sigma(y) = E \cdot \varepsilon(y) = -E \cdot y \cdot \kappa.
  \label{eq:stress_kappa}
\end{equation}

Напряжения линейно зависят от координаты $y$ и меняют знак при переходе через нейтральную ось.

\textbf{Шаг~3. Статическое соотношение (эквивалентность внутренних сил).} Внутренние силы, распределённые по сечению, должны быть эквивалентны изгибающему моменту $M(x)$ относительно нейтральной оси. Запишем условие эквивалентности:
\begin{equation}
  M(x) = \int_{A} \sigma(y) \cdot (-y) \, dA = \int_{A} E \kappa y^2 \, dA = E \kappa \int_{A} y^2 \, dA.
  \label{eq:moment_integral}
\end{equation}

Интеграл $\int_A y^2 \, dA = I_z$ является \textit{осевым моментом инерции} поперечного сечения относительно нейтральной оси. Подставляя, получаем:
\begin{equation}
  M(x) = E I_z \kappa \quad\Rightarrow\quad \kappa = \frac{M(x)}{E I_z}.
  \label{eq:kappa_M}
\end{equation}

\textbf{Шаг~4. Итоговая формула.} Подставляя $\kappa$ из~(\ref{eq:kappa_M}) обратно в~(\ref{eq:stress_kappa}), получаем формулу нормальных напряжений при изгибе:
\begin{equation}
  \sigma(x, y) = -\frac{M(x) \cdot y}{I_z}.
  \label{eq:sigma_M_y}
\end{equation}

Максимальное (по модулю) напряжение возникает в наиболее удалённых от нейтральной оси волокнах ($y = \pm y_{\max}$):
\begin{equation}
  \sigma_{\max}(x) = \frac{|M(x)| \cdot y_{\max}}{I_z} = \frac{|M(x)|}{W_z},
  \label{eq:sigma_max}
\end{equation}
\begin{conditions}
  W_z = I_z / y_{\max} & осевой момент сопротивления сечения, \textit{м}$^3$.
\end{conditions}

Условие прочности при изгибе принимает вид:
\begin{equation}
  \sigma_{\max} = \frac{|M_{\max}|}{W_z} \leq [\sigma],
  \label{eq:strength_condition}
\end{equation}
где $[\sigma]$ --- допускаемое напряжение материала~\cite{gere2018, beer2020}.

Формула~(\ref{eq:sigma_max}) реализована в платформе SopromatLab для отображения нормального напряжения в крайнем волокне при наведении курсора на эпюру изгибающего момента. Оценочное значение $y_{\max}$ вычисляется из геометрии сечения: для прямоугольника $y_{\max} = h/2$; для известного момента инерции $I$ при допущении $b \approx h$ справедливо $h = (12 I)^{1/3}$, что используется платформой в качестве первого приближения.

\subsubsection{Энергетическое обоснование принципа суперпозиции}

В рамках линейной теории упругости справедлив принцип суперпозиции: перемещения и напряжения от совокупности нагрузок равны сумме перемещений и напряжений от каждой нагрузки в отдельности. Обоснование принципа можно получить из линейности дифференциального уравнения изогнутой оси балки~(\ref{eq:euler_bernoulli}): если $v_1(x)$ и $v_2(x)$ --- решения для нагрузок $M_1(x)$ и $M_2(x)$ соответственно, то $v_1 + v_2$ является решением для $M_1 + M_2$, что непосредственно следует из линейности оператора $d^2/dx^2$~\cite{timoshenko1983}.

Принцип суперпозиции лежит в основе алгоритма построения эпюр методом суперпозиции: полная эпюра представляется как сумма эпюр от каждой отдельной нагрузки и реакции, что позволяет строить эпюры в любом сечении методом суммирования вкладов. Именно этот принцип реализован в расчётном ядре платформы SopromatLab, где значения $Q(x_k)$ и $M(x_k)$ в каждой точке дискретизации вычисляются как сумма вкладов от всех реакций и всех внешних нагрузок.
